% FY27 ONR YIP informal brief topic description. Compile using pdflatex.
% Scientific framing tailored for predrag. Program Officer line can be updated after consultation.
\documentclass[11pt,letterpaper]{article}
\usepackage[letterpaper,top=0.60in,bottom=0.58in,left=0.72in,right=0.72in]{geometry}
\usepackage{mathptmx}
\usepackage[T1]{fontenc}
\usepackage{microtype}
\usepackage{enumitem}
\usepackage{xcolor}
\usepackage[colorlinks=true,citecolor=blue,linkcolor=black,urlcolor=blue]{hyperref}
\pagestyle{empty}
\setlength{\parindent}{0pt}
\setlength{\parskip}{3pt}
\setlist[itemize]{leftmargin=1.13em,topsep=2pt,itemsep=2pt,parsep=0pt}
\emergencystretch=1.5em
\newcommand{\head}[1]{\textbf{#1}\enspace}
\newcommand{\POline}{Dr. Predrag Neskovic, Code 311 --- Mathematical Data Science}
\newcommand{\FocusArea}{Artificial Intelligence and Autonomy --- Machine Learning}
\begin{document}
\begin{center}
{\fontsize{14}{16}\selectfont\bfseries Statistical Foundations for Reliable AI Reasoning under Compute Constraints\par}
\vspace{2pt}
{\small Brief topic description for Program Officer discussion --- ONR FY27 YIP (N0001426SF004)}
\end{center}
\vspace{-2pt}
{\small \textbf{Xinya Du}, Assistant Professor of Computer Science, The University of Texas at Dallas\par
\textbf{Email:} \href{mailto:xinya.du@utdallas.edu}{xinya.du@utdallas.edu} \quad \textbf{Website:} \href{https://xinyadu.github.io}{https://xinyadu.github.io}\par
\textbf{Program Officer:} \POline\par
\textbf{Proposed ONR S\&T Focus Area:} \FocusArea\par}
\vspace{2pt}\hrule\vspace{5pt}
\head{Problem.} At inference time, language models can sample reasoning paths, score intermediate steps, or retrieve evidence. Yet majority vote and best-of-$N$ treat repeated reasoning as additional evidence, even when samples share errors. In a shipboard fuel-range estimate, fifteen of twenty solutions may reuse one misread consumption rate. Apparent consensus grows without independent confirmation. When is another reasoning trace informative, and when is a new source more valuable?
\par\head{Central hypothesis.} We hypothesize that accounting for shared errors among reasoning traces and noisy verifiers can improve accuracy at fixed cost and calibrate abstention. The mathematical questions are when step validity and error dependence are identifiable from scarce labels, and which risk-control guarantees survive dependent samples. If every trace inherits an unobserved source error, agreement alone cannot establish correctness.
\par\head{Research plan.}
\begin{itemize}
\item \textbf{Estimate information in dependent traces.} Model the validity of each reasoning step as a latent variable and each verifier score as a noisy measurement. Shared prefixes, retrieved sources, and model-specific errors create dependencies among paths. Develop structured estimators for the effective number of independent traces and the information gained from another sample, characterize identifiability with limited labels and abundant unlabeled samples, and quantify uncertainty when dependence cannot be reliably estimated.
\item \textbf{Allocate computation through sequential decisions.} Treat sampling another solution, verifying a specific step, retrieving independent evidence, and stopping or abstaining as actions with different costs. Develop value-of-information policies that trade expected error reduction against tokens, tool calls, and latency. Study selective-risk control across questions, making explicit the assumptions needed for finite-sample bounds and testing how those bounds fail under distribution shift. A stopping rule should report residual uncertainty and prefer an independent check or abstention when repeated samples offer negligible new information.
\item \textbf{Learn complementary process verifiers.} Extend synthetic-error supervision to factual, logical, and evidence-grounding checks, and measure when verifiers share rather than correct one another's mistakes. Combine their scores using the dependence model instead of ad hoc averaging. Identify regimes in which further sampling or checking adds cost but cannot resolve a common faulty premise.
\end{itemize}
\head{Evaluation and expected outcomes.} Test on MATH, MEQA, and MuSiQue, together with controlled cases where reasoning paths deliberately share prefixes or incorrect evidence. Compare against self-consistency, verifier-ranked best-of-$N$, and adaptive sampling heuristics at matched tokens and tool costs. Report accuracy--computation curves, error among accepted answers at matched coverage, calibration, and empirical validity of risk targets under dependent samples and domain shift. Ablations will test whether modeling shared sources and prefixes, rather than additional verifier calls, explains any gains.
\par\head{Naval relevance.} A shipboard planning assistant may appear confident about fuel range because many candidate calculations rely on the same rate copied from one report. Its best next step may be to verify that rate rather than generate more calculations. The research addresses Mathematical Data Science priorities in principled information integration, decision making under uncertainty, and fundamental trade-offs between inference accuracy and limited computing resources.
\par\head{Prior work and PI qualifications.} FG-PRM (\href{https://doi.org/10.18653/v1/2025.findings-emnlp.228}{\textcolor{blue}{Li et al., 2025}}, Findings of EMNLP) provides step-level error supervision; PRD (\href{https://arxiv.org/abs/2307.02762}{\textcolor{blue}{Li et al., 2024}}, TMLR) studies imperfect model judgments; and HiPRAG (\href{https://arxiv.org/abs/2510.07794}{\textcolor{blue}{Wu et al., 2026}}, ICLR) studies efficient agentic retrieval. The PI holds a 2024 NSF CAREER award and received a 2023 Amazon Research Award for process-guided fine-tuning. As a postdoc, the PI led system integration for the RESIN team's \href{https://doi.org/10.18653/v1/2022.naacl-demo.7}{\textcolor{blue}{RESIN-11}} delivery system in DARPA KAIROS.

\end{document}
